\chapter{Trigonometri pada Segitiga Siku-siku}\label{ch:g9:trig}

Pada \omterm{def:g6:shapes:triangles}{segitiga siku-siku}, begitu kamu mengetahui satu sudut lancipnya,
\emph{bentuk} segitiganya sudah tertentu --- hanya ukurannya yang dapat
berubah. Karena itu perbandingan sisinya hanya bergantung pada
sudutnya: perbandingan itulah \omterm{def:g9:trig:ratios}{kosinus}, \omterm{def:g9:trig:ratios}{sinus} dan
\omterm{def:g9:trig:ratios}{tangen}. Bersama teorema Pythagoras, ketiganya memungkinkanmu menghitung
setiap sisi dan setiap sudut \omterm{def:g6:shapes:triangles}{segitiga siku-siku} dari keterangan yang
sangat sedikit.

\section{Teorema Pythagoras, sekali lagi}

\begin{theorem}[Pythagoras]\label{thm:g9:trig:pythagoras}
Pada segitiga $ABC$ yang siku-siku di $A$, kuadrat sisi miringnya
adalah \omterm{def:g1:addition:def}{jumlah} kuadrat kedua sisi siku-sikunya:
\[
BC^2 = AB^2 + AC^2 .
\]
Sebaliknya, kalau $BC^2 = AB^2 + AC^2$ pada sebuah segitiga, maka
segitiganya siku-siku di $A$.
\end{theorem}
\admitted

\begin{example}\label{ex:g9:trig:pythagoras}
Sebuah \omterm{def:g6:shapes:triangles}{segitiga siku-siku} punya sisi siku-siku $AB = 5$ dan $AC = 12$. Maka
\[
BC^2 = 5^2 + 12^2 = 25 + 144 = 169,
\qquad\text{jadi}\quad BC = \sqrt{169} = 13 .
\]
Untuk mencari sebuah \emph{sisi siku-siku}: kalau $BC = 10$ dan
$AB = 6$, maka $AC^2 = BC^2 - AB^2 = 100 - 36 = 64$, jadi $AC = 8$.
\end{example}

\section{Kosinus, sinus, tangen}

Pada \omterm{def:g6:shapes:triangles}{segitiga siku-siku}, terhadap sudut lancip $\theta$, ketiga sisinya
punya nama: \emph{sisi miring} (yang di seberang \omterm{def:g3:shapes:rightangle}{sudut siku-sikunya},
yaitu sisi terpanjang), sisi \emph{depan} $\theta$, dan sisi
\emph{samping} $\theta$ (yaitu sisi siku-siku yang menyentuh $\theta$).

\begin{definition}[Perbandingan trigonometri]\label{def:g9:trig:ratios}
Untuk sudut lancip $\theta$ pada \omterm{def:g6:shapes:triangles}{segitiga siku-siku}:
\[
\cos\theta = \frac{\text{samping}}{\text{miring}},
\qquad
\sin\theta = \frac{\text{depan}}{\text{miring}},
\qquad
\tan\theta = \frac{\text{depan}}{\text{samping}} .
\]
Perbandingan ini\index{kosinus}\index{sinus}\index{tangen} hanya
bergantung pada $\theta$, bukan pada besar segitiganya (semua \omterm{def:g6:shapes:triangles}{segitiga
siku-siku} dengan sudut lancip yang sama adalah hasil penskalaan satu
sama lain, menurut Thales).
\end{definition}

\begin{omfigure}
\begin{tikzpicture}[scale=1.15]
  \coordinate (A) at (0,0);
  \coordinate (B) at (4.2,0);
  \coordinate (C) at (4.2,2.2);
  \draw[thick] (A) -- (B) -- (C) -- cycle;
  \draw[thin] (B) ++(-0.28,0) -- ++(0,0.28) -- ++(0.28,0);
  \draw[omThm, ->] (0.9,0) arc (0:27.7:0.9);
  \node[omThm, font=\small] at (1.25,0.21) {$\theta$};
  \node[below, font=\small, omDef] at (2.1,0) {samping};
  \node[right, font=\small, omMeth] at (4.2,1.1) {depan};
  \node[above left=-2pt, font=\small, omProp] at (2.2,1.25) {sisi miring};
  \node[below left, font=\small] at (A) {$B$};
  \node[below right, font=\small] at (B) {$A$};
  \node[above right, font=\small] at (C) {$C$};
\end{tikzpicture}

{\small Ketiga sisinya dilihat dari sudut $\theta = \widehat B$: sisi
miringnya $[BC]$, sisi sampingnya $[BA]$, sisi depannya $[AC]$.}
\end{omfigure}

\begin{proposition}[Sifat pertamanya]\label{prop:g9:trig:properties}
Untuk setiap sudut lancip $\theta$:
\begin{enumerate}
  \item $0 < \cos\theta < 1$ dan $0 < \sin\theta < 1$ (sisi siku-siku
        lebih pendek daripada sisi miringnya);
  \item $\cos^2\theta + \sin^2\theta = 1$;
  \item $\tan\theta = \dfrac{\sin\theta}{\cos\theta}$.
\end{enumerate}
\end{proposition}

\begin{proof}
Tulislah $a$ untuk sisi sampingnya, $o$ untuk sisi depannya, $h$ untuk
sisi miring sebuah \omterm{def:g6:shapes:triangles}{segitiga siku-siku} dengan sudut $\theta$.

1. Berlaku $0 < a < h$ dan $0 < o < h$, jadi kedua perbandingannya
benar-benar antara $0$ dan $1$.

2. Menurut Pythagoras, $a^2 + o^2 = h^2$. Bagilah semuanya dengan $h^2$:
\[
\cos^2\theta + \sin^2\theta
= \frac{a^2}{h^2} + \frac{o^2}{h^2}
= \frac{a^2 + o^2}{h^2} = 1 .
\]

3. $\dfrac{\sin\theta}{\cos\theta} = \dfrac{o/h}{a/h} = \dfrac oa
= \tan\theta$.
\end{proof}

\begin{method}[Mencari sebuah sisi]\label{met:g9:trig:side}
Untuk menghitung sisi yang belum diketahui ketika satu sisi dan satu
sudut lancipnya diketahui:
\begin{enumerate}
  \item tandai sudut yang diketahui lalu namai ketiga sisinya (miring,
        depan, samping) \emph{dari sudut itu};
  \item pilihlah perbandingannya (cos, sin atau tan) yang memuat sisi
        yang diketahui dan sisi yang dicari;
  \item tulislah \omterm{def:g8:equations:def}{persamaannya} lalu selesaikan;
  \item periksalah dengan akal sehat: sisi miringnya harus keluar
        sebagai yang terpanjang.
\end{enumerate}
\end{method}

\begin{example}\label{ex:g9:trig:side}
Pada sebuah \omterm{def:g6:shapes:triangles}{segitiga siku-siku}, sisi miringnya berukuran $8$ dan satu
sudutnya $35^\circ$; hitunglah sisi depannya. Perbandingan yang memuat
sisi depan dan sisi miringnya adalah sinusnya:
\[
\sin 35^\circ = \frac{\text{depan}}{8}
\quad\Longrightarrow\quad
\text{depan} = 8 \sin 35^\circ \approx 8 \times 0.574 \approx 4.6 .
\]
Untuk sisi sampingnya, pakailah kosinusnya:
$8\cos 35^\circ \approx 8 \times 0.819 \approx 6.6$. Periksalah dengan
Pythagoras: $4.6^2 + 6.6^2 \approx 21 + 43 \approx 64 = 8^2$.
\end{example}

\begin{method}[Mencari sebuah sudut]\label{met:g9:trig:angle}
Ketika dua sisinya diketahui, hitunglah perbandingan yang dibentuk
keduanya, kenalilah perbandingan itu sebagai \omterm{def:g9:trig:ratios}{kosinus}, \omterm{def:g9:trig:ratios}{sinus} atau \omterm{def:g9:trig:ratios}{tangen}
sudut yang belum diketahui, lalu terapkan fungsi \omterm{def:g8:fractions:inverse}{kebalikan} pada
kalkulatornya ($\cos^{-1}$, $\sin^{-1}$ atau $\tan^{-1}$) terhadap
perbandingan itu.
\end{method}

\begin{example}\label{ex:g9:trig:angle}
Sebuah tangga sepanjang $5$ m bersandar pada tembok, kakinya $1.4$ m
dari temboknya. Sudut $\theta$ antara tangganya dan tanah memenuhi
\[
\cos\theta = \frac{1.4}{5} = 0.28,
\qquad\text{jadi}\quad
\theta = \cos^{-1}(0.28) \approx 73.7^\circ .
\]
\end{example}

\begin{omfigure}
\begin{tikzpicture}[scale=0.85]
  \draw[thick] (0,0) -- (1.4,0);
  \draw[thick] (1.4,0) -- (1.4,4.0);
  \draw[omDef, very thick] (0,0) -- (1.4,4.0);
  \draw[thin] (1.4,0) ++(-0.25,0) -- ++(0,0.25) -- ++(0.25,0);
  \draw[omThm, ->] (0.75,0) arc (0:70.7:0.75);
  \node[omThm, font=\small] at (1.15,0.55) {$\theta$};
  \node[below, font=\small] at (0.7,0) {$1.4$ m};
  \node[omDef, font=\small, left] at (0.62,2.1) {$5$ m};
  \draw[gray, thin] (-0.7,0) -- (2.3,0);
\end{tikzpicture}

{\small Soal tangganya: sisi sampingnya ($1.4$ m) dan sisi
miringnya ($5$ m) diketahui, jadi kosinusnya memberi sudutnya.}
\end{omfigure}

\begin{example}[Sudut istimewa]\label{ex:g9:trig:special}
Dua segitiga memberi nilai yang persis. Setengah persegi bersisi 1
(yaitu \omterm{def:g6:shapes:triangles}{segitiga sama kaki} siku-siku) punya sisi miring $\sqrt2$ dan
sudut $45^\circ$:
\[
\cos 45^\circ = \sin 45^\circ = \frac{1}{\sqrt2} = \frac{\sqrt2}{2},
\qquad \tan 45^\circ = 1 .
\]
Setengah \omterm{def:g6:shapes:triangles}{segitiga sama sisi} bersisi 1 memberi sudut $30^\circ$ dan
$60^\circ$:
\[
\cos 60^\circ = \sin 30^\circ = \frac12,
\qquad
\cos 30^\circ = \sin 60^\circ = \frac{\sqrt3}{2}.
\]
\end{example}

\section{Latihan}

\begin{exercise}[$\star$]\label{exo:g9:trig:1}
Sebuah \omterm{def:g6:shapes:triangles}{segitiga siku-siku} punya sisi siku-siku $9$ dan $12$. Hitunglah
sisi miringnya. Segitiga lain punya sisi miring $17$ dan satu sisi
siku-siku $8$: hitunglah sisi siku-siku yang lain.
\end{exercise}

\begin{exercise}[$\star$]\label{exo:g9:trig:2}
Sebuah segitiga punya sisi $7$, $24$ dan $25$. Apakah ia siku-siku?
Pertanyaan yang sama untuk sisi $5$, $6$ dan $8$.
\end{exercise}

\begin{exercise}[$\star$]\label{exo:g9:trig:3}
Pada segitiga $DEF$ yang siku-siku di $D$, dengan sudut di $E$ ditulis
$\theta$: sisi yang mana yang menjadi sisi miringnya? Sisi yang mana
yang menjadi sisi depan $\theta$? Tulislah $\cos\theta$, $\sin\theta$
dan $\tan\theta$ sebagai perbandingan sisi $DE$, $DF$, $EF$.
\end{exercise}

\begin{exercise}[$\star$]\label{exo:g9:trig:4}
Pada sebuah \omterm{def:g6:shapes:triangles}{segitiga siku-siku}, sisi miringnya berukuran $10$ dan satu
sudut lancipnya $28^\circ$. Hitunglah kedua sisi siku-sikunya
(\omterm{def:g4:numbers:round}{bulatkan} ke persepuluhan;
$\cos 28^\circ \approx 0.883$, $\sin 28^\circ \approx 0.469$).
\end{exercise}

\begin{exercise}[$\star$]\label{exo:g9:trig:5}
Pada sebuah \omterm{def:g6:shapes:triangles}{segitiga siku-siku}, sisi siku-siku samping sudut $\theta$
berukuran $6$ dan sisi depannya berukuran $4.5$. Hitunglah
$\tan\theta$, lalu $\theta$
(\omterm{def:g4:numbers:round}{bulatkan} ke derajat; $\tan^{-1}(0.75) \approx 36.9^\circ$).
\end{exercise}

\begin{exercise}[$\star\star$]\label{exo:g9:trig:6}
Sebuah drone terbang pada ketinggian $120$ m. Dari seorang pengamat di
tanah, drone itu terlihat pada sudut elevasi $32^\circ$. Pada jarak
mendatar berapa dari pengamatnya drone itu berada
($\tan 32^\circ \approx 0.625$)?
\end{exercise}

\begin{exercise}[$\star\star$]\label{exo:g9:trig:7}
Sebuah jalur landai harus naik $1.2$ m dengan sudut paling besar
$5^\circ$ terhadap arah mendatar. Berapa panjang terkecil sepanjang
kemiringannya yang harus dipunyai jalur itu
($\sin 5^\circ \approx 0.0872$)? Bulatkan ke atas sampai persepuluhan
meter terdekat.
\end{exercise}

\begin{exercise}[$\star\star$]\label{exo:g9:trig:8}
Dengan memakai \cref{prop:g9:trig:properties}: sudut lancip $\theta$
memenuhi $\cos\theta = 0.6$. Hitunglah $\sin\theta$ tanpa mencari
$\theta$, lalu $\tan\theta$.
\end{exercise}

\begin{exercise}[$\star\star\star$]\label{exo:g9:trig:9}
Berilah alasan untuk nilai persis pada \cref{ex:g9:trig:special}: pada
\omterm{def:g6:shapes:triangles}{segitiga sama kaki} siku-siku dengan sisi siku-siku $1$, hitunglah sisi
miringnya; pada \omterm{def:g6:shapes:triangles}{segitiga sama sisi} bersisi $1$, hitunglah tingginya,
lalu simpulkan \omterm{def:g9:trig:ratios}{kosinus} dan \omterm{def:g9:trig:ratios}{sinus} $30^\circ$ dan $60^\circ$.
\end{exercise}

\section{Soal: Mengukur yang tidak dapat kamu jangkau}

\begin{problem}[{Soal akhir pekan --- cara dua stasiun milik juru
ukur untuk tinggi yang tidak terjangkau, kemiringan jalan, jarak ke
kaki langit, dan perlombaan tangen menuju tak berhingga}]
\label{pb:g9:trig:1}
Tidak ada meteran yang mencapai puncak gunung, puncak menara di
seberang sungai, atau kaki langit di laut --- tetapi alat ukur sudut
dan ketiga perbandingan bab ini dapat. Pusat soal ini adalah
\emph{cara dua stasiun} yang klasik milik para juru ukur, yang
menghitung sebuah tinggi tanpa pernah mendekati alasnya; di
sekelilingnya: rambu jalan, tangga, anak tangga, kelengkungan
Bumi, dan sekilas pertama sebuah fungsi yang meledak ke tak berhingga.

\medskip
\noindent\textbf{Bagian I --- Satu stasiun.}
(Nilainya: $\tan 35^\circ \approx 0.700$,
$\sin 35^\circ \approx 0.574$, $\cos 35^\circ \approx 0.819$,
$\tan 1^\circ \approx 0.0175$.)
\begin{enumerate}
  \item Dari sebuah titik $60$ m dari kaki sebuah mercusuar,
        puncaknya terlihat pada elevasi $35^\circ$ (yang diukur dari
        arah mendatar, setinggi mata). Berapa tinggi puncaknya di
        atas ketinggian mata?
  \item Sebuah rambu jalan mengumumkan kemiringan $12\,\%$: jalannya
        naik $12$ m tiap $100$ m jarak mendatar. Sudut berapa yang
        dibentuk jalannya terhadap arah mendatar? Dan berapa persen
        kemiringan yang akan dipunyai jalan bersudut $45^\circ$?
  \item Kedua sudut lancip \omterm{def:g6:shapes:triangles}{segitiga siku-siku} saling
        berpenyiku, dan sisi depan yang satu adalah sisi samping
        yang lain. Simpulkan kesamaan ``ko''-nya:
        \[
        \sin(90^\circ - \theta) = \cos\theta,
        \qquad
        \cos(90^\circ - \theta) = \sin\theta .
        \]
  \item Tangga sepanjang $8$ m bersandar pada $60^\circ$ terhadap
        tanah. Dengan memakai nilai persis pada
        \cref{ex:g9:trig:special}, sebutkan tinggi yang tercapai dan
        jarak kakinya dari temboknya --- secara persis, lalu sampai
        sentimeter.
  \item Periksalah $\cos^2\theta + \sin^2\theta = 1$ secara angka
        untuk $\theta = 35^\circ$ dan secara persis untuk
        $\theta = 60^\circ$. Mengapa kesamaan ini
        (\cref{prop:g9:trig:properties}) menjadi kebiasaan yang baik
        untuk memeriksa pekerjaan kalkulator?
\end{enumerate}

\medskip
\noindent\textbf{Bagian II --- Dua stasiun.}
Puncak sebuah gunung terlihat, tetapi kakinya tidak terjangkau. Dari
sebuah titik $A$, elevasi puncaknya $30^\circ$; sesudah berjalan
$200$ m lurus menuju gunungnya, ke $B$, elevasinya menjadi
$40^\circ$. Tulislah $h$ untuk tinggi puncaknya di atas ketinggian
mata dan $x$ untuk jarak mendatar dari $B$ ke \omterm{def:g6:lines:perp}{garis tegak lurus}
puncaknya. ($\tan 30^\circ \approx 0.577$,
$\tan 40^\circ \approx 0.839$.)
\begin{enumerate}[resume]
  \item Nyatakan $\tan 40^\circ$ dan $\tan 30^\circ$ memakai $h$,
        $x$ dan $x + 200$: yaitu dua \omterm{def:g8:equations:def}{persamaan} untuk dua yang belum
        diketahui.
  \item Dari tiap \omterm{def:g8:equations:def}{persamaannya} nyatakan jarak mendatarnya dalam
        $h$, kurangkan, lalu turunkan rumus dua stasiunnya:
        \[
        h = \frac{200}
        {\ \dfrac{1}{\tan 30^\circ} - \dfrac{1}{\tan 40^\circ}\ } .
        \]
        Hitunglah $h$ sampai puluhan meter terdekat.
  \item Hitunglah $x$ juga, lalu periksalah kedua jawabanmu dengan
        akal sehat terhadap pembidikan $30^\circ$ dari $A$.
  \item Dalam satu kalimat: mengapa cara dua stasiunnya
        sangat diperlukan --- pengukuran tunggal apa, yang mustahil
        di sini, yang akan dituntut cara satu stasiun pada
        pertanyaan 1?
  \item Sebuah menara dibidik pada $30^\circ$, lalu pada $60^\circ$
        sesudah berjalan $200$ m menujunya. Dengan memakai
        $\tan 30^\circ = \frac{1}{\sqrt3}$ dan
        $\tan 60^\circ = \sqrt3$, hitunglah tingginya
        \emph{secara persis} --- jawabannya \omterm{def:g9:arith:divisor}{kelipatan}
        $\sqrt3$ --- lalu sampai meter.
\end{enumerate}

\medskip
\noindent\textbf{Bagian III --- Ke kaki langit dan naik tembok.}
\begin{enumerate}[resume]
  \item Seberapa jauh kaki langit? Ketika berdiri dengan mata $h$
        meter di atas Bumi yang berbentuk bola berjari-jari $R$,
        garis pandangmu menyerempet bolanya: garis pandangnya,
        \omterm{def:g3:shapes:shapes}{jari-jari} ke titik serempetnya dan \omterm{def:g3:shapes:shapes}{jari-jari} ke kakimu
        membentuk \omterm{def:g6:shapes:triangles}{segitiga siku-siku}. Tunjukkan dengan Pythagoras
        (\cref{thm:g9:trig:pythagoras}) bahwa jarak $d$ ke
        kaki langit memenuhi $d^2 = 2Rh + h^2 \approx 2Rh$,
        lalu hitunglah $d$ untuk mata pada $1.80$ m
        ($R = 6.37 \times 10^6$ m).
  \item Pertanyaan yang sama dari puncak menara $324$ m. (Patokan
        kasar para pelaut: kaki langit dalam kilometer kira-kira
        $3.6\sqrt{h}$ dengan $h$ dalam meter --- periksalah kedua
        jawabanmu terhadapnya.)
  \item Ibu jarimu, yang lebarnya kira-kira $2$ cm, ditahan sejauh
        lengan, kira-kira $60$ cm dari matamu: sudut berapa yang
        ditutupinya? Bulan purnama menutupi kira-kira $0.5^\circ$:
        berapa bagian ibu jarimu yang menyembunyikan seluruh Bulan?
  \item Tangga yang nyaman punya tinggi anak tangga $17$ cm dan
        lebar pijakan $29$ cm. Sudut berapa yang didakinya?
        Bandingkan dengan $30^\circ$ pada \cref{ex:g9:trig:special},
        dan dengan tangga loteng yang curam pada $60^\circ$: tinggi
        anak tangga berapa yang \emph{diperlukannya} pada pijakan
        $29$ cm?
  \item Perlombaan \omterm{def:g9:trig:ratios}{tangennya}: hitunglah (dengan kalkulator)
        $\tan 10^\circ$, $\tan 45^\circ$, $\tan 80^\circ$,
        $\tan 89^\circ$. Apa yang terjadi ketika sudutnya mendekati
        $90^\circ$, dan mengapa (pikirkan sisi sampingnya)? Temboknya
        tegak, perbandingannya tidak punya nilai, dan grafiknya
        melesat ke tak berhingga --- itulah \emph{asimtot} pertamamu,
        makhluk yang dipelajari dengan saksama di buku Sekolah
        Menengah.

\end{enumerate}
\end{problem}
